Deux cordes sont attachées bout à bout et tendues entre deux murs avec une tension de \qty{\v{F}}{N}. La corde 1 a une longueur de \qty{\v{L1}}{m} et une masse linéique de \qty{\v{mu1}}{g/m}; la corde 2 a une longueur de \qty{\v{L2}}{m} et une masse linéique de \qty{\v{mu2}}{g/m}. Une impulsion est produite près du mur auquel est attachée la corde 1.

\begin{enumerate}[label=\alph*)]
    \item Déterminez la vitesse de l'impulsion dans chacune des cordes.
    \item Combien de temps l'impulsion transmise prend-elle pour atteindre l'extrémité de la corde 2 ?
    \item Lorsque l'impulsion arrive à la jonction, une partie est réfléchie. L'impulsion réfléchie est-elle à l'endroit ou inversée ? Justifiez.
\end{enumerate}
